\originalchapter{暴胀的退出、再加热与初始条件边界}
本章对应8.286第1--2、23讲，补充暴胀章的慢滚方程。第23讲的实际逐字稿已读完；幻灯片中Max Tegmark为Guth与Kaiser制作的WMAP曲线属于Science的例外图片，本卷不复制，也不把逐字稿误称的“Planck七年数据”沿用为图的来源。本章把模型机制、数量级例值和观测判断分开。

\originalsection{旧暴胀为何有加速却难以顺利结束}
设标量势在$\phi=0$有高于全局最低点的局部极小，场暂时陷于其间。若梯度和动能小，$p\simeq-V$、$H^2\simeq V/(3\bar M_P^2)$，可产生指数展开。经典场不能从最低的局部点自动爬过势垒；量子隧穿却可在局部形成真空泡，随后泡壁膨胀。隧穿速率来自势形和量子场论的反弹计算，不能只用势垒的高度唯一确定。

旧模型的困难是需要把能量转成足够均匀的热宇宙。若泡形成太慢，背景假真空继续迅速膨胀，泡难以并合填满空间；若太快，又难持续足够多的e折。即使泡壁碰撞可转化能量，空间强烈非均匀也未自动消失。这称平滑退出问题，指这一具体机制失败，而不是指数膨胀的微分方程本身无解。

\begin{center}\begin{tikzpicture}\begin{axis}[width=.47\textwidth,height=4.3cm,xlabel={$\phi$},ylabel={$V$},xmin=-1.25,xmax=1.25,ymin=0,ymax=.5,title={假真空与势垒}]
\addplot[structurecolor,domain=-1.25:1.25,samples=160]{(x^2-1)^2*(x^2+.2)};
\addplot[only marks,mark=*] coordinates {(0,.2)};
\end{axis}\end{tikzpicture}\hfill
\begin{tikzpicture}\begin{axis}[width=.47\textwidth,height=4.3cm,xlabel={$\phi$},xmin=-1.25,xmax=1.25,ymin=0,ymax=1.2,title={平缓山顶与连续滚动}]
\addplot[structurecolor,domain=-1.25:1.25,samples=160]{(1-x^2)^2};
\addplot[only marks,mark=*] coordinates {(0,1)};
\end{axis}\end{tikzpicture}\end{center}
图中两势是自绘多项式示意，纵横尺度任意，右图在原点的曲率是否足够小仍取决于有量纲参数。图形“较平”本身不能代替$|\eta_V|\ll1$的检查。

\originalsection{新暴胀把连续滚动与膨胀同时利用}
若势在一段区间足够平缓，场在小涨落推动后缓慢滚动，许多e折可发生在已近均匀的一个关联区域形成之后。区域被拉大，而退出通过场滚到陡区和最低点完成，避免完全依靠独立真空泡的并合。山顶是经典不稳定点，称它“暂时近真空”是在描述停留时间长，不意味着它是严格局部稳定最低点。

近山顶写$V\simeq V_f-m_t^2\phi^2/2$，$H$近常数。线性场方程为$\ddot\phi+3H\dot\phi-m_t^2\phi=0$，增长指数为
\[
 s_+=\frac{-3H+\sqrt{9H^2+4m_t^2}}2\simeq\frac{m_t^2}{3H}
 \quad(m_t^2\ll H^2).
\]
所以离开山顶时间$\sim3H/m_t^2$可远长于$H^{-1}$。这一推导具体说明“小二阶导数”何以带来长停留；仅知道$V_f$大不能保证持续时间。实际场不精确停在$\phi=0$，初始涨落和之后势形决定总e折。

\originalsection{慢滚结束后怎样产生热宇宙}
规范场在二次最低势$V\simeq m^2\phi^2/2$附近振荡。若$H\ll m$，一周期平均的动能和势能相等，故平均$w\simeq0$；在此阶段能量随$a^{-3}$稀释，而非立即变成$a^{-4}$辐射。更一般$V\propto|\phi|^p$且周期平均有效时，场的virial关系给$\langle w\rangle=(p-2)/(p+2)$。

一个简化扰动衰变模型引入速率$\Gamma_\phi$：
\[
 \dot\epsilon_\phi+3H\epsilon_\phi=-\Gamma_\phi\epsilon_\phi,
 \qquad \dot\epsilon_r+4H\epsilon_r=\Gamma_\phi\epsilon_\phi.
\]
交换项相消保证两部门总局部守恒。$H\sim\Gamma_\phi$附近若辐射已主导并快速热化，令$\epsilon_r=\pi^2g_*T_{\rm reh}^4/30$，得到
\[
 T_{\rm reh}\sim\left(\frac{90}{\pi^2g_*}\right)^{1/4}
 \sqrt{\Gamma_\phi\bar M_P}.
\]
共振预热、非扰动产生和热化延迟可能使此模型不足；最大瞬时温度也不一定等于这里定义的再加热温度。再加热不能超过可用势能预算，不能仅因课程画$10^{16}$ GeV能标便断言真实再加热温度已知。

\originalsection{课程的大统一数量级例值怎样计算}
恢复常数，用假设能标$E_G=10^{16}$ GeV作维度分析，能量密度$\epsilon_f\sim E_G^4/(\hbar c)^3$，质量等价密度$\rho_f=E_G^4/(\hbar^3c^5)$。代入$H_f=\sqrt{8\pi G\rho_f/3}$得
\[
 \rho_f\sim2.3\times10^{81}\ \mathrm{g\,cm^{-3}},\quad
 H_f^{-1}\sim2.8\times10^{-38}\ \mathrm s,\quad
 \frac c{H_f}\sim8.4\times10^{-28}\ \mathrm{cm}.
\]
第23讲一度把最后的时间写作$H_f$，随课堂提问已改为$H_f^{-1}$；本卷沿用修正后的单位。它们是\emph{假设该能标}所得样例，暴胀能标和初始时刻不能由这些数字独立测定。

从$10^{-27}$ cm扩大到1 cm需要$N\simeq\ln10^{27}\simeq62$个e折，若到10 cm约64。与原导论“约100次加倍”是同一量级，因为$100\ln2\simeq69$；具体尺寸选择不同不应冒充矛盾或精密预测。暴胀后还要接热膨胀：无净熵产生时
\[
 \frac{a_0}{a_{\rm reh}}=
 \frac{T_{\rm reh}}{T_{\gamma0}}
 \left(\frac{g_{*s,\rm reh}}{g_{*s,0}}\right)^{1/3}.
\]
不能漏掉自由度变化，更不能把$T_{\rm reh}$固定为$E_G/\kb$。这样才把一个初始因果区、e折数与今日可见区域连接起来。前面视界章的50--60 e折是指定观测模的常用估计，和任意样例区域的总e折定义并不完全相同。

\originalsection{无毛图像有明确的适用条件}
在近de Sitter的展开期，普通物质密度$a^{-3}$、辐射$a^{-4}$、空间曲率相对项$a^{-2}$均快速减小。若没有显著各向异性应力，小的均匀剪切满足$\dot\sigma+3H\sigma=0$，其能量尺度随$a^{-6}$。这说明为何许多已开始展开、真空能占优的区域趋于局部均匀和各向同性。

这种论证不证明任意初始几何必然暴胀。过大的梯度、黑洞式坍缩、特定闭合初态或不满足所用能量条件的物质都可能改变结局。正曲率切片被拉大后局部越来越平，却不因此改变其整体拓扑。de Sitter还可有不同空间切片的坐标描述，“空间平直的指数模型”与“时空完全无曲率”尤其不能混淆。

若未来始终指数展开，光从$t_i$出发的最大共动距离是$c/(a_iH)$，实际事件视界为$c/H$。它可限制外界信号对已给定de Sitter区域的影响。但有限暴胀后若接上永远减速模型，完整未来积分可能发散，从而没有永久事件视界；不能拿有限阶段的指数近似证明该区域永远隔绝任意外部变化。

\originalsection{永恒暴胀与观测成功之间还隔着什么}
局部退出不一定使\emph{全局}膨胀体积减少。作为示意，若假真空在单位时间的有效退出概率率为$\kappa$，粗略未退出体积$V_f\propto\ee^{(3H-\kappa)t}$；$3H>\kappa$可使其平均增大。这未包括泡间相关、边界和随机几何，是说明竞争的模型而非一般证明。慢滚场若一个Hubble时间的量子涨落$H/(2\pi)$足以和经典滚动$|\dot\phi|/H$竞争，也可能产生自再生区，但需检查势和随机过程。

因此永恒暴胀是某些模型及条件下的可能结果，不由“足够e折”这个条件自动推出。多真空、初始波函数、从无隧穿和景观测度均是第1--2、23讲提出的高阶方向；本卷给其物理问题和前置关系，保留完整量子宇宙学与随机测度为待补范围。

CMB的近标度不变、近绝热扰动以及峰结构与一类暴胀初始条件相容。位置与相对高度还依赖重子、暗物质、几何和散射史；归一振幅往往需拟合势的参数。课程展示的六参数拟合不能说成所有点均为无参数预测，也不能由平均谱与角功率谱的成功选出唯一暴胀势或证明多重宇宙。真实数据与模型对应关系，须和CMB章、热史章及距离拟合一同理解。
